\chapter{Virtual Memory and Page Replacement}\label{ch:vm}
\chapterguide{Understand contiguous allocation from Chapter~\ref{ch:allocation}.}{Explain pages and frames, translate a virtual page through a page table, define page faults, and trace FIFO, LRU, and LFU replacement.}

\section{Why virtual memory changes the problem}
Paging divides virtual memory into fixed-size \term{pages} and physical memory into equal-size \term{frames}. A page table records which frame currently contains each resident virtual page. Because pages need not occupy adjacent frames, paging removes external fragmentation from physical allocation, though the last page of an object can still create internal fragmentation.

\section{Address translation}
Divide the virtual address by the page size. The quotient is the page number;
the remainder is the offset. The page table replaces the page number with a
frame number, while the offset stays unchanged. In compact form:
\texttt{physical address = frame $\times$ page size + offset}.

\begin{mlfigure}
\centering
\begin{tikzpicture}[
  cell/.style={draw=MLBlue,fill=MLPaleBlue,minimum width=1.6cm,minimum height=0.7cm,
               font=\small\color{MLNavy},align=center},
  off/.style={draw=MLOrange,fill=MLPaleOrange,minimum width=1.6cm,minimum height=0.7cm,
              font=\small\color{MLNavy},align=center},
  pte/.style={draw=MLTeal,fill=MLPaleTeal,minimum width=1.2cm,minimum height=0.58cm,
              font=\scriptsize\color{MLNavy}},
  lab/.style={font=\scriptsize\color{MLMuted},align=center},
  ar/.style={-{Stealth[length=2.2mm]},thick,draw=MLTeal}
]
% The two address fields are STACKED, not side by side, so the page-number
% path and the offset path never cross each other or the page table.

% ---- headings ----
\node[lab] at (0,3.5)   {\textbf{Virtual address}\\ $v = 9\,000$, \ $P = 4\,096$};
\node[lab] at (4.9,4.0) {\textbf{Page table}};
\node[lab] at (9.6,3.5) {\textbf{Physical address}};

% ---- page table, positioned so entry 2 lines up with the page row ----
\node[pte] (p0) at (4.9,3.2) {0 $\to$ 7};
\node[pte] (p1) at (4.9,2.6) {1 $\to$ 3};
\node[pte,draw=MLRed,fill=MLPaleRed,very thick] (p2) at (4.9,2.0) {2 $\to$ 5};
\node[pte] (p3) at (4.9,1.4) {3 $\to$ 1};

% ---- page row (y = 2.0) and offset row (y = 0.5) ----
\node[cell] (vp) at (0,2.0) {page 2};
\node[off]  (vo) at (0,0.5) {offset 808};
\node[cell] (pf) at (9.6,2.0) {frame 5};
\node[off]  (po) at (9.6,0.5) {offset 808};

% ---- arithmetic, below its own column ----
\node[lab] at (0,-0.35)   {$\lfloor 9000/4096 \rfloor = 2$\\ $9000 \bmod 4096 = 808$};
\node[lab] at (9.6,-0.35) {$5 \times 4096 + 808$\\ $= 21\,288$};

% ---- translated: page number in, frame out ----
\draw[ar] (vp.east) -- node[lab,above]{look up} (p2.west);
\draw[ar] (p2.east) -- node[lab,above]{gives frame} (pf.west);

% ---- untranslated: the offset passes straight under the table ----
\draw[ar,draw=MLOrange] (vo.east) --
      node[lab,above]{copied unchanged --- it never enters the page table} (po.west);
\end{tikzpicture}

\smallskip
{\scriptsize Only the \emph{page number} is translated. The offset passes through
untouched, which is why the page size fixes how many bits each part uses --- and
why page size is always a power of two.}
\end{mlfigure}

\begin{keyidea}
Translation replaces the page number and copies the offset. If a virtual address is contiguous inside one page, its physical address is contiguous too; across a page boundary it need not be, because the next page may live in any frame.
\end{keyidea}

\section{Page faults}
A \term{page fault} occurs when a referenced virtual page is not currently resident in a usable physical frame. The OS must locate the page on backing storage (or otherwise materialize it), obtain a frame, update mapping metadata, and restart/resume the faulting instruction as appropriate.

\section{Replacement policies}
When no free frame exists, the OS needs a victim page.

\begin{mltable}
\begin{tabularx}{\linewidth}{@{}>{\bfseries\leavevmode\color{MLNavy}}p{0.16\linewidth}p{0.29\linewidth}X@{}}
\toprule
\rowcolor{MLPaleBlue!92}
Policy & Victim rule & State required \\
\midrule
FIFO & Page resident for the longest time & Arrival/load order \\
LRU & Page whose most recent use is farthest in the past & Recency information \\
LFU & Page with the smallest reference count & Frequency counts plus a tie-breaker \\
\bottomrule
\end{tabularx}
\end{mltable}

\section{FIFO}
FIFO treats loaded pages like a queue. It is easy to implement but does not use actual access behavior. FIFO can exhibit \term{Belady's anomaly}: for some reference strings, adding more frames can increase the number of faults.

\section{LRU}
LRU approximates the principle of temporal locality: a page used recently is likely to be used again soon. Exact LRU can be expensive, so real systems often use hardware-assisted approximations.

\section{LFU}
LFU retains frequently referenced pages. Pure cumulative counts can preserve stale historical popularity, so practical frequency-based schemes often age counts or use bounded windows.

\begin{workedexample}
With three frames and reference string 1, 2, 3, 1, 4: the first three references fault and fill the frames. The second reference to 1 is a hit. On reference 4, FIFO evicts 1 if 1 was loaded earliest, whereas LRU evicts 2 because 2 is the least recently used at that moment.
\end{workedexample}

\section{Locality explains why replacement policies can work}
Programs often reuse recently accessed instructions and data (temporal locality) and access addresses close to one another (spatial locality). Page-replacement policies try to exploit these patterns indirectly. LRU uses recent history; LFU uses frequency history; FIFO deliberately ignores reference history and uses load order only.

\begin{keyidea}
A page-replacement algorithm is a prediction policy: it tries to choose a resident page whose removal will hurt future execution as little as possible. The policy never knows the future perfectly, so the workload matters.
\end{keyidea}

\section{How to trace replacement professionally}
A trustworthy trace shows the reference, frame contents after the reference, hit/fault result, and enough policy state to justify a victim. Totals are useful only after individual decisions can be checked.

\begin{quickcheck}
\begin{enumerate}
  \item Which part of a virtual address remains unchanged during page-table translation?
  \item What is the difference between a page hit and a page fault?
  \item Why can LRU and LFU choose different victims even on the same reference string?
\end{enumerate}
\end{quickcheck}

\labsection{8}{Virtual memory and page replacement}{sec:lab08}
\begin{labproject}{Trace FIFO, LRU, and LFU using the same reference string}

\oslabflow{Chapter 8: page hits, page faults, replacement state.}{Run three policies on one reference string so the victim rule is the only major change.}{Justify at least one replacement decision from each policy's stored state.}

\section*{Source code}
\ossourcecard{\texttt{fifo.c}}{Replace the page that has been resident the longest.}{\texttt{next} rotates through frame positions; page-use history is ignored.}
\lstinputlisting[style=oslabc]{all_experiments/Lab08/fifo.c}

\ossourcecard{\texttt{lru.c}}{Replace the least recently used resident page.}{\texttt{last\_used[]} stores the most recent reference step for each frame.}
\lstinputlisting[style=oslabc]{all_experiments/Lab08/lru.c}

\ossourcecard{\texttt{lfu.c}}{Replace the least frequently used resident page; use recency to break ties.}{Two pieces of policy state are tracked: \texttt{frequency[]} and \texttt{last\_used[]}.}
\lstinputlisting[style=oslabc]{all_experiments/Lab08/lfu.c}

\section*{Task 1 --- Compile and collect traces}
\begin{lstlisting}[style=osbash]
cd all_experiments/Lab08
for f in fifo lru lfu; do
  gcc -Wall -Wextra -std=c11 "$f.c" -o "$f"
  ./$f
  echo
 done
\end{lstlisting}
Fill in the final hit/fault totals.

\begin{mltable}
\renewcommand{\arraystretch}{1.18}
\begin{tabularx}{\linewidth}{@{}
  >{\raggedright\arraybackslash}p{0.20\linewidth}
  *{2}{>{\raggedright\arraybackslash}X}@{}}
\toprule
\rowcolor{UTPPaleBlue!92}
\textbf{Policy} & \textbf{Hits} & \textbf{Faults} \\
\midrule
FIFO & & \\
LRU & & \\
LFU & & \\
\bottomrule
\end{tabularx}
\end{mltable}

\section*{Task 2 --- Explain one victim from each policy}
Choose a step where all three frames are full and a fault occurs. For FIFO identify the oldest frame position; for LRU compare last-use steps; for LFU compare frequencies and then the tie-break rule if needed.

\section*{Task 3 --- Verify the beginning by hand}
Manually trace the first six references: 7, 0, 1, 2, 0, 3. Compare your frame table with all three programs. The first few decisions should make the policies' state differences visible.

\section*{Task 4 --- Controlled edit}
Choose one policy. In Vim change the final reference in its array from 2 to 7. Predict whether that final access is a hit or a fault for that policy, recompile, and verify.

\begin{quickcheck}
\begin{enumerate}
  \item What event causes a page fault?
  \item What state does LRU need that FIFO does not?
  \item Why can two replacement policies choose different victims for the same reference?
\end{enumerate}
\end{quickcheck}

\begin{labdeliverable}
Submit the three totals, one justified victim decision per policy, your six-reference hand trace, and the prediction/result for the final-reference edit.
\end{labdeliverable}
\end{labproject}

